\chapter{Addison's Disease}
\index{Addison's Disease}

\section*{Definition and Overview}
Addison's disease, also called primary adrenal insufficiency, is a rare disorder in which the adrenal glands, two small glands that sit on top of the kidneys, do not produce enough of the hormones cortisol and aldosterone. Cortisol helps the body respond to stress, regulate blood sugar, and fight inflammation, while aldosterone controls salt and water balance and blood pressure.

Without treatment, Addison's disease is life-threatening, but with lifelong hormone replacement most people lead normal lives. Because the condition develops slowly, it is often diagnosed only after months of vague symptoms, and it is frequently missed until an acute crisis forces urgent attention.

\section*{Epidemiology}
Addison's disease is rare, affecting roughly 1 in 10,000 people. It can occur at any age but is most often diagnosed in adults between 30 and 50 years, and it is slightly more common in women than in men.

In developed countries the most common cause is autoimmune destruction of the adrenal glands. The condition is more likely in people who have other autoimmune conditions, such as thyroid disease, type 1 diabetes, or vitiligo.

\section*{Aetiology and Pathophysiology}
In primary adrenal insufficiency the adrenal glands themselves fail. In about 8 in 10 cases the cause is autoimmune disease, in which the body's immune system mistakenly attacks and gradually destroys the outer layer (cortex) of the glands.

Less common causes include tuberculosis, fungal infections, cancer that has spread to the glands, and, rarely, bleeding into the glands. When both cortisol and aldosterone are missing, the body cannot regulate stress, blood sugar, blood pressure, or salt balance properly, which produces the characteristic symptoms. Because cortisol also provides feedback that suppresses the pituitary's production of ACTH, its loss allows ACTH to rise, which stimulates the darkening of skin seen in primary disease.

\section*{Clinical Features}
The disease develops slowly, and symptoms may be vague for months:

\begin{itemize}
    \item \textbf{Fatigue:} progressive tiredness and weakness out of proportion to activity
    \item \textbf{Weight loss:} loss of appetite and unintentional weight loss
    \item \textbf{Low blood pressure:} dizziness or light-headedness on standing
    \item \textbf{Skin darkening:} pigmentation in scars, skin creases, and gums (a feature of primary disease)
    \item \textbf{Salt craving:} a strong desire for salty foods
    \item \textbf{Gastrointestinal:} abdominal pain, nausea, and diarrhoea
    \item \textbf{Biochemical changes:} low blood sugar, low sodium, and high potassium
    \item \textbf{Other features:} loss of body hair and low mood
\end{itemize}

\section*{Differential Diagnosis}
Because the early symptoms are vague, Addison's disease is often mistaken for other conditions:

\begin{itemize}
    \item Depression and chronic fatigue
    \item Thyroid disease
    \item Anaemia or chronic infection
    \item Diabetes, especially with weight loss and fatigue
    \item Gastrointestinal disorders causing weight loss and vomiting
    \item Anorexia nervosa
\end{itemize}

\section*{When to Seek Medical Care}
Anyone with persistent fatigue, weight loss, dizziness, and darkening skin should see a doctor for assessment. Once diagnosed, a person should also seek medical care for any illness, injury, surgery, or stress that could trigger an adrenal crisis, so that steroid doses can be adjusted in time.

\textbf{Contact your local emergency services for:}
\begin{itemize}
    \item Severe vomiting or diarrhoea that prevents keeping fluids down
    \item Extreme weakness, collapse, or loss of consciousness
    \item Very low blood pressure or shock
    \item Confusion, agitation, or severe abdominal pain
    \item Any illness in a person with Addison's disease who cannot take their medicines
\end{itemize}

\section*{Diagnosis and Investigations}
\begin{itemize}
    \item \textbf{Blood tests:} low sodium and high potassium, and a low morning cortisol level
    \item \textbf{ACTH measurement:} the hormone that stimulates the adrenal glands is typically elevated in primary disease
    \item \textbf{Short Synacthen test:} the definitive test, measuring cortisol before and after an injection of synthetic ACTH
    \item \textbf{Aldosterone and renin:} to assess the salt-regulating system
    \item \textbf{Imaging:} a CT scan of the adrenal glands to look for structural causes
    \item \textbf{Other tests:} screening for coexisting autoimmune conditions and adrenal antibodies
\end{itemize}

\section*{Management}
\begin{itemize}
    \item \textbf{Hormone replacement:} daily cortisol replacement (hydrocortisone) and aldosterone replacement (fludrocortisone)
    \item \textbf{Sick day rules:} follow the person's individualized written plan for adjusting steroid replacement during illness, injury, or stress
    \item \textbf{Emergency injection:} a hydrocortisone injection kit for use in an adrenal crisis
    \item \textbf{Medical alert:} wearing a bracelet or carrying a card stating the condition and medicines
    \item \textbf{Regular monitoring:} blood tests, blood pressure checks, and dose adjustment
    \item \textbf{Review of other medicines:} some medicines, such as certain antifungal drugs, interact with steroid replacement
\end{itemize}

\section*{Complications}
\begin{itemize}
    \item Adrenal crisis, a life-threatening emergency with very low blood pressure and shock
    \item Low blood sugar (hypoglycaemia)
    \item Osteoporosis from long-term steroid use
    \item Weight gain, high blood pressure, and diabetes from excessive steroid doses
    \item Other autoimmune diseases, which may develop over time
    \item Depression and reduced quality of life
\end{itemize}

\section*{Prognosis and Outcomes}
With correct treatment and careful management, most people with Addison's disease can expect a near-normal life expectancy and a good quality of life. The main danger is an adrenal crisis, which occurs most often when steroid doses are not increased during illness or are missed.

People who learn to manage their condition well, follow sick day rules, and carry an emergency kit generally do very well. Education of the patient and their family is the single most effective safeguard against serious complications.

\section*{Prevention and Self-Care}
\begin{itemize}
    \item Take replacement steroids every day, exactly as prescribed
    \item Learn and follow the individualized sick-day plan; contact an urgent medical service if vomiting or diarrhoea prevents tablets being absorbed
    \item Carry an emergency hydrocortisone kit and a medical alert card
    \item Tell all doctors and dentists about the condition before any procedure
    \item Never stop steroids suddenly, as this can trigger a crisis
    \item Keep regular appointments for monitoring and dose review
    \item Seek early medical help for any infection or injury
\end{itemize}

\section*{Sources and Further Reading}
The following organisations provide reliable information on Addison's disease and adrenal insufficiency.

\begin{itemize}
    \item NHS. \textit{Information on Addison's disease, diagnosis, and treatment.} https://www.nhs.uk/conditions/
    \item Mayo Clinic. \textit{Guide to Addison's disease and its management.} https://www.mayoclinic.org/
    \item MedlinePlus. \textit{Health information on adrenal gland disorders.} https://medlineplus.gov/
    \item MSD Manual. \textit{Professional information on adrenal insufficiency.} https://www.msdmanuals.com/
    \item NICE. \textit{Clinical guidance relevant to endocrine disorders and steroid treatment.} https://www.nice.org.uk/
    \item Endocrine Society. \textit{Clinical practice guideline: primary adrenal insufficiency.} https://www.endocrine.org/clinical-practice-guidelines/primary-adrenal-insufficiency
\end{itemize}
